\lecture{5}
\title{Bayesian Linear Regression}

\begin{document}

\subsection{Problems with Linear Regression}

Recall that $\hat{\theta} = (X^{\top}X)^{-1}X^{\top}Y$
is the minimizer of
$\mathrm{RSS}(\theta) := (X\theta - Y)^{\top}(X\theta - Y)$
(the ``Residual Sum of Squares'').
One possible problem is that $X$ might not be full rank;
two or more features might be perfectly collinear,
making $\hat{\theta}$ non-unique.

\begin{center}
    \frame{\includegraphics[width=0.4\textwidth]{lecture-05/not-full-rank}}
\end{center}

Maybe in practice, there's some tiny amount of noise
in our observations, so $X$ is ``almost always'' guaranteed to be full rank.
That doesn't actually fix the issue at hand, though---the effect of that noise
is still to effectively pick one of the infinitely-many maximum-likelihood
choices for $\hat{\theta}$ at \emph{random}.

\begin{remark}
    To clarify, collinearity of features isn't uncommon at all;
    take ``credit card limit'' and ``credit card rating''
    as a practical example of two nearly-collinear features.
\end{remark}

The issue is that linear regression is \textbf{overconfident};
it always returns a \emph{single} point estimate $\hat{\theta}$.
Rather than decisively estimating the exact value of $\hat{\theta}$,
perhaps we should take a \textbf{Bayesian approach}:
let incoming data $\{(x^{(n)}, y^{(n)})\}_{n = 1}^N$
\emph{update our priors} on the distribution of $\theta$ instead.

\subsection{Bayesian Linear Regression: Theory}

Here's the model;
refer to \href{/18.650/lectures/21}{Lecture 21} of 18.650
for a reminder on Bayesian inference.
\begin{itemize}
    \item Our prior on $\theta$ is that
    $\theta \sim \mathcal{N}(\mu_0, \Sigma_0)$,
    for arbitrary initial constants $\mu_0$ and $\Sigma_0$.
    \item We will assume that data $(x^{(n)}, y^{(n)})$
    is such that $y^{(n)} \sim \mathcal{N}(\theta^{\top}x^{(n)}, \sigma^2)$
    for all $n$.
\end{itemize}
A prior and a model for our data
are enough for us to determine the posteriors
that result from reading in data.

\begin{theorem}{Linear Regression Posterior}
    Suppose we collect $N$ data points,
    represented by a matrix $X$ and a column vector $Y$.
    Then the posterior is
    $[\theta \mid X, Y] \sim \mathcal{N}(\mu_N, \Sigma_N)$,
    where $\mu_N$ and $\Sigma_N$ satisfy:
    \[
    \Sigma_N^{-1} = \Sigma_0^{-1} + \frac{1}{\sigma^2}X^{\top}X
    ~~~ \text{ and } ~~~
    \Sigma_N^{-1}\mu_N = \Sigma_0^{-1}\mu_0 + \frac{1}{\sigma^2}X^{\top}Y
    \]
\end{theorem}

\begin{remark}
    Goes without saying, but\dots
    even though these aren't explicit formulas for $\Sigma_N$ and $\mu_N$,
    they do provide enough information
    to fully compute them if desired.
\end{remark}

\begin{proof}
    Plain computation, with some simplifying tricks.
    Recall that $p(\theta \mid X, Y) \propto p(\theta) \cdot L_N(\theta)$,
    where $p(\theta)$ is the PDF of a multivariable Gaussian,
    and $L_N(\theta)$ is a product of $N$ single-variable Gaussian PDFs.
    \[
        p(\theta) = \frac{1}{\sqrt{(2\pi)^D \det(\Sigma_0)}}
        \exp\left(-\frac{1}{2}(\theta - \mu_0)^{\top}\Sigma_0^{-1}(\theta - \mu_0)\right)
        ~~~ \text{ and } ~~~
        L_N(\theta) = \frac{1}{\sqrt{(2\pi\sigma^2)^N}}
        \exp\left(-\frac{1}{2\sigma^2}\|Y - X\theta\|^2\right).
    \]
    See \href{/18.650/lectures/6/#multivariate-gaussian}{Lecture 6}
    of 18.650 for a derivation of the multivariable case.

    For ease of computation,
    it helps to instead work with $\log p(\theta \mid X, Y)$.
    Also recall that we only care about
    $p(\theta \mid X, Y)$ up to a proportionality constant
    (as indicated by the $\propto$ symbol),
    so after taking logs,
    we can throw away any \textcolor{teal}{$+C$} terms.

    We'll compute the logarithm in two steps,
    starting with $\log p(\theta)$.
    (Note that all terms in $\log p(\theta)$ are real numbers.)

    \begin{align*}
        \log p(\theta)
        & = -\frac{1}{2}(\theta - \mu_0)^{\top}\Sigma_0^{-1}(\theta - \mu_0)
        \textcolor{teal}{+C} \\
        & = -\frac{1}{2}\left[\theta^{\top}\Sigma_0^{-1}\theta
        - \mu_0^{\top}\Sigma_0^{-1}\theta
        - \theta^{\top}\Sigma_0^{-1}\mu_0
        \textcolor{teal}{+ \mu_0^{\top}\Sigma_0^{-1}\mu_0}\right]
        \textcolor{teal}{+C} \\
        & = -\frac{1}{2}\left[\theta^{\top}\Sigma_0^{-1}\theta
        \textcolor{purple}{- \mu_0^{\top}\Sigma_0^{-1}\theta}
        \textcolor{purple}{- \theta^{\top}\Sigma_0^{-1}\mu_0}\right]
        \textcolor{teal}{+C} & (\bigstar_1) \\
        & = \textcolor{blue}{\theta^{\top}\left(
            -\frac{1}{2}\Sigma_0^{-1}
        \right)\theta}
        \textcolor{purple}{+ \left(\mu_0^{\top}\Sigma_0^{-1}\right)\theta}
        \textcolor{teal}{+C} & (\bigstar_2)
    \end{align*}
    Two subtle simplifying tricks just took place:
    \begin{itemize}
        \item To arrive at line $(\bigstar_1)$,
        we absorbed the $\textcolor{teal}{+ \mu_0^{\top}\Sigma_0^{-1}\mu_0}$
        into the $\textcolor{teal}{+C}$
        because it is a real-number constant (i.e. has no $\theta$s).
        \item To arrive at line $(\bigstar_2)$,
        we replaced \textcolor{purple}{$\theta^{\top}\Sigma_0^{-1}\mu_0$}
        with \textcolor{purple}{$\mu_0^{\top}\Sigma_0^{-1}\theta$}.
        We can do this because
        they are real numbers and transposes of each other,
        so they are equal.
        In particular, $(\Sigma_0^{-1})^{\top} = \Sigma_0^{-1}$
        because covariance matrices $\Sigma_0$ are symmetric.
    \end{itemize}
    Using some similar tricks,
    we can compute $\log L_N(\theta)$ as well.
    \begin{align*}
        \log L_N(\theta)
        & = -\frac{1}{2\sigma^2}(Y - X\theta)^{\top}(Y - X\theta)
        \textcolor{teal}{+C} \\
        & = -\frac{1}{2\sigma^2}\left[
            \textcolor{teal}{Y^{\top}Y}
            \textcolor{purple}{-\theta^{\top}X^{\top}Y}
            \textcolor{purple}{-Y^{\top}X\theta}
            + \theta^{\top}X^{\top}X\theta
        \right] \textcolor{teal}{+C} \\
        & = \textcolor{blue}{
            \theta^{\top}\left(
                -\frac{1}{2\sigma^2} X^{\top}X
            \right)\theta
        } \textcolor{purple}{
            +\left(\frac{1}{\sigma^2}Y^{\top}X\right)\theta
        } \textcolor{teal}{+C}.
    \end{align*}
    So altogether, we have the following:
    \[
        (\clubsuit_1) ~~~
        \log p(\theta \mid X, Y)
        = \textcolor{blue}{
            \theta^{\top}\left(
                -\frac{1}{2}\left[
                    \Sigma_0^{-1} + \frac{1}{\sigma^2}X^{\top}X
                \right]
            \right) \theta
        } \textcolor{purple}{+ \left(
            \mu_0^{\top}\Sigma_0^{-1} + \frac{1}{\sigma^2}Y^{\top}X
        \right)\theta} \textcolor{teal}{+C}.
    \]
    We just need to compare this with
    the log of the PDF of $\mathcal{N}(\mu_N, \Sigma_N)$\dots
    well, we already computed what the logarithm looks like---it's
    just line $(\bigstar_2)$ from earlier!
    \[
        (\clubsuit_2) ~~~
        \log \left[\text{PDF of } \mathcal{N}(\mu_N, \Sigma_N)\right]
        = \textcolor{blue}{\theta^{\top}\left(
            -\frac{1}{2}\Sigma_N^{-1}
        \right)\theta}
        \textcolor{purple}{+ \left(\mu_N^{\top}\Sigma_N^{-1}\right)\theta}
        \textcolor{teal}{+C}
    \]
    Pattern-matching $(\clubsuit_1)$ against $(\clubsuit_2)$
    (and taking the transpose of the latter equation\dots)
    yields the desired result. ~ $\blacksquare$
\end{proof}

\begin{remark}
    The constant simplification of everything
    into the form $\theta^{\top}A\theta + b\theta + c$
    isn't a coincidence---every multivariable quadratic polynomial
    is expressible in such a form!  See
    \href{/18.701/lectures/16/#foreshadowing-quadratic-polynomials-are-bilinear-forms}{Lecture 16}
    of 18.701 for further detail.
\end{remark}

\subsection{Bayesian Linear Regression: Performance}

The whole point of Bayesian linear regression
is to avoid overconfidence.
To what extent have we been successful?

\begin{quote}
    \textbf{Example Set-Up.}
    We draw features $x^{(n)} \sim \mathrm{Unif}[-1, 1]$
    and generate responses via
    $y^{(n)} \sim \mathcal{N}(\theta_1^*x^{(n)} + \theta_0^*, \sigma^2)$,
    where $\theta_0^* = -0.3$, $\theta_1^* = 0.5$, and $\sigma = 0.2$.
    The algorithm assumes the prior $\theta \sim \mathcal{N}(0, I_2)$.
\end{quote}

Under this set-up,
Bayesian linear regression performs expectedly well;
as more data comes in, the posterior becomes
more focused around the correct value of $\theta$.

\begin{center}
    \frame{\includegraphics[width=0.4\textwidth]{lecture-05/example-1}}
\end{center}

But ordinary linear regression would do just as well here;
there is no feature collinearity here.

\begin{quote}
    \textbf{Altered Set-Up.}
    Everything held constant, except the only feature value
    that is observed is $x^{(n)} = 1$.
\end{quote}

This is equivalent to feature collinearity.
(Recall that, for \emph{affine} linear regression,
all feature vectors have an extra ``1'' appended to their front
to allow for ``y-intercepts''; see
\href{/18.650/lectures/24/#affine-linear-regression}{Lecture 24}
of 18.650.)
Bayesian linear regression still performs well here!

\begin{center}
    \frame{\includegraphics[width=0.4\textwidth]{lecture-05/example-2}}
\end{center}

The regression is able to be confident about $\theta^*$
in some parameter directions, but not all parameter directions.
There's still a one-dimensional space of
potential values of $\theta^* = (\theta_0^*, \theta_1^*)$
that the regression sees as possible---still
a strong improvement over the two-dimensional prior!

\begin{quote}
    \textbf{Further Altered Set-Up.}
    The Bayesian regression believes responses are generated via
    $y^{(n)} \sim \mathcal{N}(\theta_1 x^{(n)} + \theta_0, \sigma^2)$,
    but they're actually generated by
    $y^{(n)} \sim \mathcal{N}(0.5 - (x^{(n)})^2, \sigma^2)$.
\end{quote}

Bayesian regression can't save you from model \emph{misspecification}.
As more data arrives, the posterior becomes more and more resolute
in its estimate for $\theta$,
even though the regression hasn't come close
to actually understanding the response generation.

\begin{center}
    \frame{\includegraphics[width=0.4\textwidth]{lecture-05/example-3}}
\end{center}

\end{document}